% !TeX program = xelatex
%==============================================================
%  Python金融数据分析 · 第3章 NumPy数组计算
%  基于《Python金融大数据分析（第2版）》第4章内容改编
%  授课对象：本科金融系学生（已学完第2章Python基础语法）
%  编译：xelatex chapter3.tex（编译两遍）
%==============================================================
\documentclass[9pt]{ctexbeamer}

%%%%-----导入宏包-----%%%%
\usepackage{style/shubeamer}   % SHU 主题（配色、帧标题、页脚、Python代码样式）
\usepackage{amsmath,amsfonts,amssymb,bm}
\usepackage{graphicx,hyperref,url}
\usepackage{subcaption}
\usepackage{booktabs}
\usepackage[super,square]{natbib}
\usepackage{wrapfig}          % 文字绕排
\usepackage{calligra}         % 结尾页花体字

%%%%-----设置字体-----%%%%
\setmainfont{CMU Serif}
\setsansfont{CMU Sans Serif}

%%%%-----常用自定义颜色（正文可直接使用）-----%%%%
\definecolor{nvidia}{RGB}{102,156,28}
\definecolor{red}{RGB}{184,13,73}
\definecolor{orange}{RGB}{242,151,36}
\definecolor{darkteal}{RGB}{43,106,108}
\definecolor{darkblue}{RGB}{4,37,58}

%%%%-----每节开始时自动显示当前节目录-----%%%%
\AtBeginSection[]
{
	\begin{frame}<beamer>
		\frametitle{\textbf{目录}}
		\textbf{\tableofcontents[currentsection]}
	\end{frame}
}

%%%%-----设定图片存放路径-----%%%%
\graphicspath{{figures/}}

%%%%-----首页信息设置-----%%%%
\AtBeginDocument{%
	%%%%----短标题[显示在页脚]{封面标题}
	\title[Python金融数据分析 · 第3章]{\fontsize{16pt}{22pt}\selectfont {Python金融数据分析}}
	%%%%----副标题（本章主题）
	\subtitle{\fontsize{9pt}{14pt}\selectfont \textit{第3章\quad NumPy数组计算}}
	%%%%----个人信息
	\author[王伟冠, 吕文俊]{%
		\begin{tabular}{ll}%
			主讲人： & 王伟冠, 吕文俊\tabularnewline%
			院\quad 系： & 经济学院金融系\tabularnewline
		\end{tabular}
	}
	%%%%----机构信息
	\institute[SHU]{上海大学}
	%%%%----日期信息
	\date[\today]{\today}}

%%%%-----数学宏（向量、矩阵、算子等，见 style/macros.tex）-----%%%%
\input{style/macros.tex}


\begin{document}

%% 封面
\begin{frame}
	\titlepage
\end{frame}

%% 目录页
\section*{目录}
\begin{frame}
	\frametitle{\textbf{目录}}
	\textbf{\tableofcontents}
\end{frame}


%==============================================================
\section{为什么需要NumPy？}
%==============================================================

\begin{frame}{本章学习目标}
	\begin{itemize}
		\item 理解为什么金融数据分析离不开 NumPy
		\item 掌握 \texttt{ndarray} 的创建、属性与索引切片
		\item 掌握\textbf{向量化运算}：一行代码对整个数组做计算
		\item 掌握\textbf{布尔索引}：按条件筛选数据
		\item 掌握常用统计函数（均值、标准差、累计等）
		\item 能够用 NumPy 计算收益率序列与波动率等风险指标
	\end{itemize}
	\vspace{0.3em}
	\begin{block}{前置知识}
		第2章的列表、循环、函数——本章会把它们"升级"为数组写法。
	\end{block}
\end{frame}

\begin{frame}[fragile]{回顾：上一章的做法}
	\begin{lstlisting}
prices = [100.0, 102.0, 101.0, 105.0]
returns = []
for i in range(1, len(prices)):
    r = (prices[i] - prices[i-1]) / prices[i-1]
    returns.append(r)
\end{lstlisting}
	\begin{itemize}
		\item 只有 4 个价格没问题；但真实数据动辄\alert{数千个交易日、成百上千只股票}。
		\item Python 的 \texttt{for} 循环每一步都要检查类型、逐个搬运，\textbf{速度慢}；
		\item 代码冗长，意图（"算收益率"）被循环细节淹没。
	\end{itemize}
	\begin{alertblock}{本章的目标}
		用 NumPy 把 4 行循环压缩成 1 行，并且快上百倍。
	\end{alertblock}
\end{frame}

\begin{frame}[fragile]{NumPy是什么？}
	\begin{itemize}
		\item \textbf{NumPy}（Numerical Python）是Python科学计算的基石。
		\item 核心对象：\texttt{ndarray}——存储\textbf{同一类型}数据的多维数组。
		\item 底层用 C 实现，运算直接作用于整块内存，速度远超Python循环。
		\item pandas、scikit-learn 等后续工具都构建在 NumPy 之上。
	\end{itemize}
	\begin{lstlisting}
import numpy as np     # 约定俗成的导入方式

prices = np.array([100.0, 102.0, 101.0, 105.0])
returns = prices[1:] / prices[:-1] - 1    # 一行搞定！
\end{lstlisting}
	\begin{exampleblock}{约定}
		\texttt{import numpy as np} 是全世界通用的写法，请照抄。
	\end{exampleblock}
\end{frame}


%==============================================================
\section{创建数组}
%==============================================================

\begin{frame}[fragile]{从列表创建与基本属性}
	\begin{lstlisting}
a = np.array([100.0, 102.0, 101.0, 105.0])
a            # array([100., 102., 101., 105.])

type(a)      # numpy.ndarray
a.shape      # (4,)      形状：长度为4的一维数组
a.dtype      # float64   元素类型：全数组统一
len(a)       # 4
\end{lstlisting}
	\begin{itemize}
		\item \texttt{np.array(列表)}：把Python列表转换为数组。
		\item \textbf{类型统一}：数组中所有元素同一类型（如全是 \texttt{float64}），这正是快的原因。
		\item \texttt{shape} 描述"形状"，是后面理解二维数据的关键。
	\end{itemize}
\end{frame}

\begin{frame}[fragile]{快速生成常用数组}
	\begin{lstlisting}
np.zeros(5)           # array([0., 0., 0., 0., 0.])
np.ones(3)            # array([1., 1., 1.])

np.arange(0, 10, 2)   # array([0, 2, 4, 6, 8])  类似range
np.linspace(0, 1, 5)  # array([0., 0.25, 0.5, 0.75, 1.])
                      # 0到1之间均匀取5个点（含端点）
\end{lstlisting}
	\begin{itemize}
		\item \texttt{zeros/ones}：初始化占位数组（如预存结果）。
		\item \texttt{arange(起, 止, 步长)}：整数序列，与 \texttt{range} 类似但返回数组。
		\item \texttt{linspace(起, 止, 个数)}：画利率区间、画函数图像时常用。
	\end{itemize}
	\begin{exampleblock}{金融场景}
		给一组折现率 \texttt{np.linspace(0.01, 0.10, 10)}，批量观察债券价格变化。
	\end{exampleblock}
\end{frame}

\begin{frame}[fragile]{二维数组：数据的"表格"形态}
	\begin{lstlisting}
# 3天 x 2只股票 的收盘价：每行一天，每列一只股票
p = np.array([[25.0, 100.0],
              [25.5,  99.0],
              [26.0, 102.0]])

p.shape      # (3, 2)   3行2列
p.ndim       # 2        维度数
p.size       # 6        元素总数
\end{lstlisting}
	\begin{itemize}
		\item 列表中套列表 $\rightarrow$ 二维数组，可以看作\textbf{表格/矩阵}。
		\item \texttt{shape} 变为 \texttt{(行数, 列数)}。
		\item 金融数据的典型形态：\alert{行=日期，列=资产}。
	\end{itemize}
\end{frame}

% 自编教学补充；阅读来源见本章 README。
\begin{frame}[fragile]{2.3A 用成绩表理解 shape 与 axis}
\small 每行是一名学生，每列是一门课。axis=0 沿学生方向汇总，得到各科均分；axis=1 沿科目方向汇总，得到各学生均分。
\medskip
\begin{lstlisting}
scores = np.array([[80, 90, 70], [60, 75, 90]])
print(scores.shape)        # 2名学生、3门课
print(scores.mean(axis=0)) # [70. 82.5 80.]
print(scores.mean(axis=1)) # [80. 75.]
\end{lstlisting}
\end{frame}

\begin{frame}[fragile]{形状变换与拼接}
	\begin{lstlisting}
a = np.arange(6)             # [0 1 2 3 4 5]
a.reshape(2, 3)              # 变成2行3列
a.reshape(-1, 2)             # -1 表示"自动算"：3行2列

np.array([1, 2]) + 0         # 数组拼接用专门函数：
np.concatenate([np.array([1, 2]), np.array([3, 4])])
                             # array([1, 2, 3, 4])
\end{lstlisting}
	\begin{alertblock}{易错点}
		数组的 \texttt{+} 是\textbf{逐元素相加}，不是拼接！拼接请用 \texttt{np.concatenate}。\\
		\texttt{np.array([1,2]) + np.array([3,4])} 得到 \texttt{[4, 6]}。
	\end{alertblock}
\end{frame}


%==============================================================
\section{索引与切片}
%==============================================================

\begin{frame}[fragile]{一维数组：与列表几乎相同}
	\begin{lstlisting}
prices = np.array([30.1, 30.5, 29.8, 31.2, 30.9])

prices[0]        # 30.1    索引从0开始
prices[-1]       # 30.9    倒数第一个
prices[1:3]      # array([30.5, 29.8])   含头不含尾
prices[1:]       # 除第一个外的所有元素
prices[:-1]      # 除最后一个外的所有元素
\end{lstlisting}
	\begin{itemize}
		\item 索引、切片规则与列表完全一致——旧知识直接复用。
		\item \texttt{prices[1:]} 与 \texttt{prices[:-1]} 这对组合后面算收益率要用，务必熟练。
	\end{itemize}
	\begin{alertblock}{注意}
		数组切片返回的是\textbf{视图}（与原数组共享内存）：修改切片会影响原数组，与列表不同。
	\end{alertblock}
\end{frame}

\begin{frame}[fragile]{二维数组：先行后列}
	\begin{lstlisting}
p = np.array([[25.0, 100.0],
              [25.5,  99.0],
              [26.0, 102.0]])

p[0, 0]      # 25.0   第0行第0列
p[1, 1]      # 99.0   第1行第1列
p[0]         # 第0行 -> array([25., 100.])
p[:, 1]      # 所有行、第1列 -> 股票2的价格序列
p[1:, :]     # 第1行到末尾的所有行
\end{lstlisting}
	\begin{itemize}
		\item 写法：\texttt{数组[行, 列]}；冒号 \texttt{:} 表示"全部"。
		\item \texttt{p[:, 1]} 即"取出某只股票的整个价格序列"，金融分析中的高频操作。
	\end{itemize}
\end{frame}


%==============================================================
% 自编教学补充；阅读来源见本章 README。
\begin{frame}[fragile]{3.3 切片视图与显式复制}
\small 基本切片通常共享原数组的数据。需要独立修改时加 copy()。这与 Python 列表的切片不同。
\medskip
\begin{lstlisting}
measurements = np.array([10, 20, 30, 40])
view = measurements[1:3]
independent = measurements[1:3].copy()
view[0] = 99
print(measurements)  # [10 99 30 40]
print(independent)   # [20 30]
\end{lstlisting}
\end{frame}
\begin{frame}[fragile]{课堂练习：3.3 切片视图与显式复制}
改成 measurements[[1, 2]] 后再修改结果，原数组会变化吗？
\medskip
\begin{block}{检查思路}
先写出预期结果，再运行。参考答案见配套 Notebook 的折叠单元格。
\end{block}
\end{frame}

\section{向量化运算}
%==============================================================

\begin{frame}[fragile]{数组与标量：整体运算}
	\begin{lstlisting}
prices = np.array([100.0, 102.0, 101.0, 105.0])

prices + 1       # 每个元素加1
prices * 1.05    # 每个元素涨5%
prices ** 2      # 每个元素平方
1 / prices       # 每个元素取倒数
\end{lstlisting}
	\begin{itemize}
		\item \textbf{向量化}：运算自动作用于数组中的\alert{每一个}元素，无需循环。
		\item 运算不改变原数组，而是返回\textbf{新数组}。
		\item 读代码时把数组想象成"一整列Excel数据同时参与计算"。
	\end{itemize}
\end{frame}

\begin{frame}[fragile]{数组与数组：逐元素配对}
	\begin{lstlisting}
a = np.array([1.0, 2.0, 3.0])
b = np.array([10., 20., 30.])

a + b        # array([11., 22., 33.])
a * b        # array([10., 40., 90.])
b / a        # array([10., 10., 10.])
\end{lstlisting}
	\begin{itemize}
		\item 两个数组运算：\textbf{位置一一对应}（第0个对第0个、第1个对第1个……）。
		\item 这正对应金融场景：
		\begin{itemize}
			\item "今日价 / 昨日价"：两个错开一位的价格序列相除；
			\item "组合收益"：持仓权重数组 $\times$ 个股收益数组。
		\end{itemize}
	\end{itemize}
	\begin{alertblock}{形状要匹配}
		参与运算的两个数组形状必须一致（或满足广播规则），否则报错。
	\end{alertblock}
\end{frame}

% 自编教学补充；阅读来源见本章 README。
\begin{frame}[fragile]{4.2A 广播：每科加不同的分数}
\small (2, 3) 与 (3,) 可以逐列配对。广播从末尾维度比较，长度相同或其中一个为 1 才能兼容。
\medskip
\begin{lstlisting}
scores = np.array([[80, 90, 70], [60, 75, 90]])
bonus = np.array([2, 0, 5])
adjusted = scores + bonus
print(adjusted)
student_mean = scores.mean(axis=1, keepdims=True)
print(student_mean.shape)  # (2, 1)
print(scores - student_mean)
\end{lstlisting}
\end{frame}
\begin{frame}[fragile]{课堂练习：4.2A 广播：每科加不同的分数}
将每一科减去该科均分；解释结果为什么每列均值为 0。
\medskip
\begin{block}{检查思路}
先写出预期结果，再运行。参考答案见配套 Notebook 的折叠单元格。
\end{block}
\end{frame}

\begin{frame}[fragile]{例：一行代码计算收益率}
	\begin{lstlisting}
prices = np.array([100.0, 102.0, 101.0, 105.0])

# 简单收益率：(今日-昨日)/昨日
ret = prices[1:] / prices[:-1] - 1
ret        # array([0.02, -0.0098, 0.0396])

# 对数收益率（后文常用）
logret = np.log(prices[1:] / prices[:-1])
\end{lstlisting}
	\begin{itemize}
		\item \texttt{prices[1:]}：今日价序列；\texttt{prices[:-1]}：昨日价序列；相除即逐日配对。
		\item 上一章的4行循环，现在\alert{1行}；数据越大，优势越明显。
		\item \texttt{np.log()} 是向量化函数，对整个数组取自然对数。
	\end{itemize}
\end{frame}

\begin{frame}[fragile]{快多少？速度对比}
	\begin{lstlisting}
big = np.arange(1_000_000)          # 0到999999

%timeit [x ** 2 for x in big]       # Python循环：约200毫秒
%timeit big ** 2                    # NumPy向量化：约1毫秒
\end{lstlisting}
	\begin{itemize}
		\item \texttt{\%timeit} 是 Notebook 计时魔法命令（下节 notebook 中可亲自体验）。
		\item 同样的一百万元素平方运算，NumPy 快约 \alert{100倍以上}。
		\item 原因：类型统一 + 底层C实现 + 无需逐步类型检查。
	\end{itemize}
	\begin{block}{经验法则}
		看到"对每个元素做同样的事"，先想向量化，再考虑循环。
	\end{block}
\end{frame}

\begin{frame}[fragile]{常用向量化数学函数}
	\begin{table}
		\centering
		\begin{tabular}{lll}
			\toprule
			函数 & 作用 & 金融用途举例 \\
			\midrule
			\texttt{np.log(x)} / \texttt{np.exp(x)} & 对数 / 指数 & 对数收益率、连续复利 \\
			\texttt{np.sqrt(x)} & 平方根 & 波动率换算 \\
			\texttt{np.abs(x)} & 绝对值 & 偏差幅度 \\
			\texttt{np.maximum(a,b)} & 逐元素取大 & 期权收益 $\max(S-K,0)$ \\
			\texttt{np.round(x, n)} & 四舍五入 & 报表展示 \\
			\texttt{np.cumsum(x)} & 累计和 & 累计收益、净值曲线 \\
			\texttt{np.diff(x)} & 相邻差分 & 价格变动额 \\
			\bottomrule
		\end{tabular}
	\end{table}
	\begin{exampleblock}{期权到期收益}
		\texttt{np.maximum(S - K, 0)}：行权价 $K$，到期价序列 $S$，一行算出所有情形的收益。
	\end{exampleblock}
\end{frame}


%==============================================================
\section{布尔索引：按条件筛选}
%==============================================================

\begin{frame}[fragile]{比较运算返回布尔数组}
	\begin{lstlisting}
ret = np.array([0.02, -0.01, 0.04, -0.03, 0.01])

ret > 0      # array([True, False, True, False, True])
\end{lstlisting}
	\begin{itemize}
		\item 数组的比较运算也是\textbf{向量化}的：逐元素判断，返回等长的布尔数组。
		\item 布尔数组可以当"筛选器"使用（下一页）。
		\item 组合条件要用 \texttt{\&}（且）、\texttt{|}（或），且每个条件加括号：
	\end{itemize}
	\begin{lstlisting}
(ret > -0.01) & (ret < 0.01)     # 小幅波动日
\end{lstlisting}
	\begin{alertblock}{易错点}
		数组条件中 \texttt{and}、\texttt{or} 会报错，必须用 \texttt{\&}、\texttt{|}。
	\end{alertblock}
\end{frame}

\begin{frame}[fragile]{用布尔数组筛选与统计}
	\begin{lstlisting}
ret = np.array([0.02, -0.01, 0.04, -0.03, 0.01])

ret[ret > 0]          # array([0.02, 0.04, 0.01]) 只留上涨日
ret[ret < -0.02]      # array([-0.03])            大跌日

np.sum(ret > 0)       # 3     True按1计数：上涨天数
(ret > 0).mean()      # 0.6   上涨日占比
\end{lstlisting}
	\begin{itemize}
		\item \texttt{数组[布尔数组]}：只保留对应位置为 \texttt{True} 的元素。
		\item \texttt{np.sum(布尔数组)}：统计满足条件的个数——非常实用的技巧。
		\item 金融应用：统计涨停天数、回撤超过阈值的次数、极端波动日。
	\end{itemize}
\end{frame}

\begin{frame}[fragile]{np.where：条件赋值}
	\begin{lstlisting}
ret = np.array([0.02, -0.01, 0.04])

np.where(ret > 0, 1, 0)
# array([1, 0, 1])   条件为真取1，否则取0

np.where(ret > 0, "涨", "跌")
# array(['涨', '跌', '涨'])
\end{lstlisting}
	\begin{itemize}
		\item \texttt{np.where(条件, 真时的值, 假时的值)}：向量化的 \texttt{if/else}。
		\item 应用：构造交易信号——收益为正记1（持有多头），为负记0。
		\item 这比写 \texttt{for} 循环套 \texttt{if} 简洁且快得多。
	\end{itemize}
\end{frame}


%==============================================================
\section{统计函数}
%==============================================================

\begin{frame}[fragile]{常用聚合函数}
	\begin{lstlisting}
ret = np.array([0.02, -0.01, 0.04, -0.03, 0.01])

ret.sum()        # 0.03     求和
ret.mean()       # 0.006    均值（日均收益率）
ret.std()        # 0.026... 标准差（波动性）
ret.max()        # 0.04     最大值
ret.min()        # -0.03    最小值
ret.argmax()     # 2        最大值的"位置"
\end{lstlisting}
	\begin{itemize}
		\item "把一串数压缩成一个数"的函数叫\textbf{聚合函数}。
		\item 写法两种等价：\texttt{ret.mean()} 或 \texttt{np.mean(ret)}。
		\item \textbf{标准差}是金融中风险（波动率）的度量，务必熟悉。
	\end{itemize}
\end{frame}

\begin{frame}[fragile]{二维数组与 axis 参数}
	\begin{lstlisting}
# 3天 x 2只股票的日收益率
r = np.array([[ 0.02, -0.01],
              [-0.01,  0.03],
              [ 0.04,  0.01]])

r.mean()          # 0.0133...  全部元素的均值
r.mean(axis=0)    # 按"列"求均值 -> 每只股票的日均收益
r.mean(axis=1)    # 按"行"求均值 -> 每天的组合平均收益
\end{lstlisting}
	\begin{itemize}
		\item \texttt{axis=0}：沿行的方向压缩 $\rightarrow$ 结果对应每一\textbf{列}；
		\item \texttt{axis=1}：沿列的方向压缩 $\rightarrow$ 结果对应每一\textbf{行}。
		\item 记忆口诀：\alert{axis 是被"压掉"的那一维}。
	\end{itemize}
\end{frame}

\begin{frame}[fragile]{累计与差分}
	\begin{lstlisting}
prices = np.array([100., 102., 101., 105.])

np.diff(prices)     # array([2., -1., 4.]) 每日涨跌"额"
np.cumsum([1, 1, 1, 1])   # array([1, 2, 3, 4]) 累计和
\end{lstlisting}
	\begin{itemize}
		\item \texttt{np.diff}：相邻元素相减，长度少1——快速得到每日变动。
		\item \texttt{np.cumsum}：逐步累加——从"每日现金流"得到"累计现金流"，\\
			从"对数收益率"恢复"累计对数收益"。
	\end{itemize}
	\begin{exampleblock}{小例子}
		对数收益率序列取 \texttt{np.cumsum}，再取 \texttt{np.exp}，即得净值曲线（相对期初价格）。
	\end{exampleblock}
\end{frame}


%==============================================================
% 自编教学补充；阅读来源见本章 README。
\begin{frame}[fragile]{6.4 非金融小案例：左右走动}
\small 每步向左或向右移动一格，起点为 0。先生成步长，再用 cumsum 得到位置；随机种子使课堂输出可复现。
\medskip
\begin{lstlisting}
walk_rng = np.random.default_rng(7)
directions = walk_rng.integers(0, 2, size=12)
steps = np.where(directions == 0, -1, 1)
positions = np.r_[0, steps.cumsum()]
print(steps)
print(positions)
print("最远距离：", np.abs(positions).max())
\end{lstlisting}
\end{frame}
\begin{frame}[fragile]{课堂练习：6.4 非金融小案例：左右走动}
改成 100 步，并计算走动后位于起点右侧的比例；分母是否包括起点？
\medskip
\begin{block}{检查思路}
先写出预期结果，再运行。参考答案见配套 Notebook 的折叠单元格。
\end{block}
\end{frame}

\section{综合案例与小结}
%==============================================================

\begin{frame}[fragile]{综合案例：模拟股价并测风险}
	\begin{lstlisting}
rng = np.random.default_rng(8)          # 随机数生成器(固定种子)
n = 250                                  # 250个交易日

# 模拟日收益率：均值0.04%，日波动1.5%
rets = 0.0004 + 0.015 * rng.standard_normal(n)

# 由收益率还原价格曲线（期初100元）
prices = 100 * np.exp(np.cumsum(rets))

print(f"日均收益率：{rets.mean():.4%}")
print(f"日波动率：  {rets.std():.4%}")
print(f"年化波动率：{rets.std()*np.sqrt(252):.2%}")
print(f"最高价/最低价：{prices.max():.2f} / {prices.min():.2f}")
\end{lstlisting}
	\begin{itemize}
		\item \textbf{年化波动率 = 日标准差 $\times\sqrt{252}$}，业界标准做法。
		\item 用到的全部是本章工具：向量化、\texttt{cumsum}、\texttt{exp}、聚合函数。
	\end{itemize}
\end{frame}

\begin{frame}[fragile]{续：进一步分析}
	\begin{lstlisting}
# 上涨日占比
print(f"上涨日占比：{(rets > 0).mean():.2%}")

# 单日跌幅超过3%的"极端日"
extreme = rets[rets < -0.03]
print(f"极端下跌天数：{len(extreme)}")

# 最大回撤的粗略观察：历史最高价 vs 当前价
peak = np.maximum.accumulate(prices)
drawdown = prices / peak - 1
print(f"最大回撤：{drawdown.min():.2%}")
\end{lstlisting}
	\begin{itemize}
		\item \texttt{np.maximum.accumulate}：到每个时点为止的历史最高价。
		\item \textbf{最大回撤}是基金与策略最常用的风险指标之一。
		\item 整段分析没有写一个 \texttt{for} 循环。
	\end{itemize}
\end{frame}

\begin{frame}{本章小结}
	\begin{itemize}
		\item \textbf{创建}：\texttt{np.array}、\texttt{zeros/ones}、\texttt{arange/linspace}；属性 \texttt{shape/dtype}。
		\item \textbf{索引切片}：与列表同规则；二维写 \texttt{[行, 列]}，\texttt{:} 表示全部。
		\item \textbf{向量化}：数组运算逐元素进行，一行顶循环，速度快百倍；\\
		\alert{数组的 \texttt{+} 是相加不是拼接}。
		\item \textbf{布尔索引}：\texttt{a[a>0]} 筛选；\texttt{np.sum(条件)} 计数；条件组合用 \texttt{\&}、\texttt{|}。
		\item \textbf{统计}：\texttt{mean/std/max/min}；\texttt{axis} 是被压掉的维度。
		\item \textbf{金融工具箱}：收益率 \texttt{p[1:]/p[:-1]-1}、年化波动率 $\sigma\times\sqrt{252}$、\texttt{cumsum/diff}、\texttt{np.where} 信号。
	\end{itemize}
\end{frame}

\begin{frame}{课后任务与下章预告}
	\begin{block}{课后任务}
		\begin{enumerate}
			\item 运行课堂 notebook，修改参数观察输出变化。
			\item 给定价格数组 \texttt{np.array([50, 52, 51, 55, 53, 56])}：\\
			计算日收益率、找出最大单日涨幅、统计上涨天数。
			\item 选做：用 \texttt{np.where} 构造信号"收益为正买1手、为负卖1手"，\\
			并计算该策略的每日收益（信号$\times$次日收益，注意错位）。
		\end{enumerate}
	\end{block}
	\begin{exampleblock}{下章预告}
		数组只装数字，真实行情还有日期、股票代码、公司名。\\
		下一章学习 \textbf{pandas}：带标签的表格数据，金融数据分析的主力工具。
	\end{exampleblock}
	\begin{center}
		\vspace{0.5em}
		本章内容改编自教材第4章\citep{hilpisch2018python}。
	\end{center}
\end{frame}


%==============================================================
%  附录：参考文献 + 结尾页
%==============================================================
\begin{frame}{扩展阅读与练习来源}
\begin{itemize}
\item Wes McKinney, \textit{Python for Data Analysis}, 3rd ed.
\item \url{https://wesmckinney.com/book/numpy-basics}
\item 新增中文说明、合成数据与练习为本课程自行编写。
\item 阅读正文理解概念；在配套 Notebook 中运行、修改、核对。
\end{itemize}
\end{frame}

\appendix

\begin{frame}[allowframebreaks]{参考文献}
	\bibliographystyle{style/gbt7714-2005}
	\bibliography{reference.bib}
\end{frame}

\begin{frame}
	\begin{center}
		{\Huge\calligra Thanks for your attention!}
		\vspace{1cm}

		{\Huge Q \& A}
	\end{center}
\end{frame}

\end{document}
